\documentclass[12pt]{article}
\usepackage[utf8]{inputenc}
\usepackage[T1]{fontenc}
\usepackage{hyperref}
\title{Yotti Language Resources for K-12 Curriculum Development}
\author{Compiled by Claude Sonnet 5.5 (claude-sonnet-5-5) for Andr\'as Kornai}
\date{}
\begin{document}
\maketitle

\section{Introduction}
Yotti (also spelled Yoti; ISO 639-3 \texttt{yot}; Glottocode yott1234) is a small Atlantic-Congo language of Taraba State, northeastern Nigeria. Wikipedia gives about 3,000 speakers (2011). Joshua Project reports a population of about 4,400 and lists Wonkorok as an alternate name. Glottolog lists the language as not endangered. Searches found almost no pedagogical or descriptive material.

\section{Primary Educational Resources}
No primers, textbooks, or school materials in Yotti turned up in search results.

\section{Linguistic and Reference Materials}
\begin{itemize}
\item Wikipedia, ``Yoti language'': \url{https://en.wikipedia.org/wiki/Yoti_language}. It mentions a Paradisec collection of Roger Blench's materials that includes Yoti. I did not locate the collection itself.
\item Glottolog entry for Yotti: \url{https://glottolog.org/resource/languoid/iso/yot}
\item ASJP word-list entry: \url{https://asjp.clld.org/languages/YOTI}
\end{itemize}

\section{Religious and Community Resources}
\begin{itemize}
\item Joshua Project, Yotti in Nigeria (demographic and people-group profile): \url{https://joshuaproject.net/people_groups/20549}
\item Joshua Project language page: \url{https://joshuaproject.net/index.php/languages/yot}
\end{itemize}
No Yotti scripture or community organization was found.

\section{Digital and Audio Resources}
No audio, video, or software resources were found beyond the archival mention above.

\section{Curriculum Development Context}
Very little exists for Yotti. A curriculum would have to start with new documentation: an orthography, a lexicon, and recordings. The Paradisec/Blench materials are the first thing to follow up. Nothing found indicates that the language is taught in schools.
\end{document}